% !TEX root = IE3_expE_prelab.tex
\documentclass[lualatex,a4paper,10pt]{jlreq}
% Shared LuaLaTeX preamble.
\usepackage{luatexja-fontspec}
\setmainfont{TeX Gyre Termes}
\setsansfont{TeX Gyre Heros}
\setmonofont{Latin Modern Mono}
\setmainjfont[BoldFont=HaranoAjiGothic-Medium]{HaranoAjiMincho-Regular}
\setsansjfont[BoldFont=HaranoAjiGothic-Bold]{HaranoAjiGothic-Medium}
\usepackage[top=22mm,bottom=22mm,left=23mm,right=23mm]{geometry}
\usepackage{amsmath,amssymb,booktabs,array,longtable,tabularx}
\usepackage{graphicx,xcolor,tikz,circuitikz}
\usetikzlibrary{arrows.meta,positioning,calc}
\usepackage{enumitem,fancyhdr,listings,needspace}
\usepackage[unicode,colorlinks=true,linkcolor=labblue,urlcolor=labteal]{hyperref}
\usepackage{bookmark}
\definecolor{labblue}{RGB}{30,70,112}
\definecolor{labteal}{RGB}{0,100,105}
\definecolor{labgray}{gray}{0.95}
\ModifyHeading{section}{font={\large\sffamily\bfseries\color{labblue}}}
\ModifyHeading{subsection}{font={\normalsize\sffamily\bfseries\color{labteal}}}
\setlength{\parindent}{1\zw}
\setlength{\parskip}{3pt}
\setlength{\emergencystretch}{2em}
\renewcommand{\arraystretch}{1.35}
\setlist{itemsep=3pt,topsep=4pt,leftmargin=2\zw}
\newcolumntype{Y}{>{\raggedright\arraybackslash}X}
\newcolumntype{C}[1]{>{\centering\arraybackslash}p{#1}}
\pagestyle{fancy}
\fancyhf{}
\fancyhead[L]{\small\color{labblue} IE3 後期・電子工学実験}
\fancyhead[R]{\small テーマ\LabCode}
\fancyfoot[C]{\thepage}
\setlength{\headheight}{16pt}
\DeclareOldFontCommand{\rm}{\normalfont\rmfamily}{\mathrm}
\lstset{basicstyle=\small\ttfamily,columns=fullflexible,keepspaces=true,
 breaklines=true,frame=single,backgroundcolor=\color{labgray},
 numbers=left,numberstyle=\tiny,showstringspaces=false,aboveskip=8pt,belowskip=8pt}
\newcommand{\blank}{\rule{22mm}{0.3pt}}
\newcommand{\answerline}{\par\noindent\rule{\linewidth}{0.25pt}\par}
\newcommand{\writing}[1]{\par\Needspace{\dimexpr#1+5mm\relax}\noindent%
 \fcolorbox{labblue!35}{white}{\parbox[c][#1][t]{\dimexpr\linewidth-2\fboxsep-2\fboxrule\relax}{\vspace{2mm}\noindent\strut\hfill\par}}\par}
\newcommand{\hint}[1]{\par\smallskip\noindent{\sffamily\bfseries\color{labteal}記述の視点：}#1\par}
\newcommand{\note}[1]{\par\smallskip\noindent\fcolorbox{labblue!45}{labblue!5}{%
 \parbox{\dimexpr\linewidth-2\fboxsep-2\fboxrule\relax}{#1}}\par\smallskip}
\newcommand{\eqerror}{\varepsilon=\frac{x_{\mathrm{meas}}-x_{\mathrm{th}}}{x_{\mathrm{th}}}\times100\ \%}
\newcommand{\LabSource}{徳山工業高等専門学校 情報電子工学科，
 『2026年度 電子工学実験（3年後期）テキスト』，テーマ\LabCode，
 \expandafter\path\expandafter{\SourcePdf}。}
\newcommand{\reporttitle}{
 \noindent{\small\sffamily\color{labteal}実験報告書・記入ガイド}\par
 \noindent{\LARGE\sffamily\bfseries \LabCode\quad\LabTitle}\par\smallskip
 \noindent 氏名：\blank\quad 班：\blank\quad 実験日：\blank\par
 \note{目的・原理・手順を確認し、実測値と自分の考察を記入する。
 考察のガイドは、検討する事象・使う測定結果・記述の方針を示す。}
}
\newcommand{\prelabtitle}{
 \noindent{\small\sffamily\color{labteal}事前学習・前提知識プリント}\par
 \noindent{\LARGE\sffamily\bfseries \LabCode\quad\LabTitle}\par\smallskip
 \noindent 氏名：\blank\quad 班：\blank\quad 予習日：\blank\par
 \note{用語、回路の動き、計算、測定表への記入の順に読む。
 例題の値は説明用であり、実験結果ではない。}
}
\newcommand{\tocrow}[3]{\hyperref[#1]{#2}& #3 &\pageref*{#1}\\[2mm]}
\newcommand{\documentmap}[1]{
 \section*{目次と概要}
 \noindent 青色の項目をクリックすると該当箇所へ移動する。\par
 {\small\begin{longtable}{@{}p{0.29\linewidth}p{0.61\linewidth}r@{}}
 \toprule {\color{labblue}\bfseries 項目} & {\color{labblue}\bfseries 読むと分かること} & 頁\\\midrule
 \endhead #1\bottomrule
 \end{longtable}}
 \clearpage
}
\newenvironment{terms}{\par\smallskip\begin{longtable}{@{}p{0.23\linewidth}p{0.73\linewidth}@{}}
 \toprule {\color{labteal}\bfseries 用語}&{\color{labteal}\bfseries 意味・回路での役割}\\\midrule\endhead}
 {\bottomrule\end{longtable}}
\newcommand{\term}[2]{\textbf{#1}&#2\\[2mm]}
\newcommand{\guidepoint}[4]{
 \Needspace{32mm}\subsection{#1}
 \noindent{\bfseries\color{labteal}考える事象：}#2\par
 \noindent{\bfseries\color{labteal}参照する結果：}#3\par
 \noindent{\bfseries\color{labteal}記述の方針：}#4\par
}
\newcommand{\reportclosing}[1]{
 \clearpage\section{結論のまとめ方}\label{sec:closing}
 \hint{#1}
 \ConclusionGuide
 \begin{enumerate}
 \item 目的に対応する最も重要な実測結果を、測定条件と数値を添えて述べる。
 \item 原理と一致した点・一致しなかった点を示す。原因は確認済みの事実と推測を分ける。
 \item 実施範囲で言えることと、未確認の条件を一文ずつまとめる。
 \end{enumerate}
 \note{書き出しの型：「○○の条件で△△を測定したところ、□□となった。
 これは◇◇と整合する。ただし××は確認しておらず、一般化には追加測定が必要である。」
 空欄は自分の実測値と根拠で埋める。}
 \noindent 結論（実験後に記入）：\writing{30mm}
 \noindent 改善案・次に確認したい条件：\writing{16mm}
 \section*{参考文献}\label{sec:references}
 \begin{enumerate}
 \item \LabSource
 \item 使用したデータシート・書籍・ウェブ資料（著者、題名、該当箇所、URL、参照日）を追加する。
 \end{enumerate}
}
\newcommand{\prelabclosing}{
 \section*{準備・出典}\label{sec:prelabsource}
 \begin{itemize}[nosep]
 \item 資料、筆記具、記録用紙、電卓と、実験条件・機器設定の記録欄。
 \item 確認問題の解答と、分からない用語・回路点への具体的な質問。
 \end{itemize}
 {\small\noindent\LabSource\par}
}
\newcommand{\simpleflow}[1]{
 \begin{center}
 \begin{tikzpicture}[node distance=5mm and 5mm,
 every node/.style={font=\small},box/.style={draw=labblue,fill=labblue!4,rounded corners=1mm,
 align=center,minimum height=10mm,text width=30mm}]
 #1
 \end{tikzpicture}
 \end{center}
}

\newcommand{\LabCode}{E}
\newcommand{\LabTitle}{マルチバイブレータ}
\newcommand{\SourcePdf}{IE3_expE_A4.pdf}
\begin{document}
\prelabtitle
\documentmap{
\tocrow{sec:auto1}{1 用語を一つずつ理解する}{言葉の意味、単位、回路や測定での役割を確認する。}
\tocrow{sec:auto2}{2 充放電がスイッチを切り替える}{この項目の仕組みと実験で確認すべき点を整理する。}
\tocrow{sec:auto3}{3 RC充放電と状態}{指数変化としきい値から持続時間を読む。}
\tocrow{sec:auto4}{4 555の時間計算}{充放電経路から幅・周期・デューティを計算する。}
\tocrow{sec:auto5}{5 離散回路を読む}{部品と信号の経路を追い、状態が変化する理由を理解する。}
\tocrow{sec:auto6}{6 測定で区別する量}{周期とパルス幅、入力と出力を区別する。}
\tocrow{sec:auto7}{7 測定からレポートへ：時間と状態を対応させる}{実測値から計算・作図・記述へ進む順序を例題で練習する。}
\tocrow{sec:auto8}{8 確認問題}{自分で計算・説明して、実験に必要な理解を確かめる。}
\tocrow{sec:auto9}{解答の目安}{数値と理由を照合し、単位や論理の取り違えを見直す。}
\tocrow{sec:auto10}{補足資料}{メーカー公式資料で方式・端子・動作条件を確認する。}
\tocrow{sec:prelabsource}{準備・出典}{持参物と資料を確認し、具体的な質問を準備する。}
}

\section{用語を一つずつ理解する}\label{sec:auto1}
\begin{terms}
\term{マルチバイブレータ}{状態を切り替えて矩形波や一定幅のパルスを作る回路。安定状態の数で非安定・単安定・双安定に分類する。}
\term{非安定・単安定}{非安定は外部トリガなしに二状態を交互に繰り返す。単安定は待機状態からトリガで一時的に変化し、一定時間後に戻る。}
\term{時定数RC}{コンデンサ電圧の変化の速さを決める時間。RCの単位は秒。RC秒で必ずしきい値に達するという意味ではない。}
\term{比較器としきい値}{二電圧の大小を判定する回路と、その判定の境目。555は電源の約1/3、2/3を境目に状態を切り替える。}
\term{ラッチ}{状態を記憶する回路。555の内部では比較器の判定で切り替わり、出力と放電トランジスタを制御する。}
\term{トリガ・微分}{動作を開始させる信号がトリガ。RC微分回路は急な変化を短いパルスに変え、切替え時刻を取り出すために使う。}
\term{周期・幅・デューティ}{周期は同じエッジから次の同じエッジまで。幅は一つの状態が続く時間。デューティは高レベル時間の周期に対する割合である。}
\end{terms}
\section{充放電がスイッチを切り替える}\label{sec:auto2}
\subsection{離散トランジスタ非安定回路}
片方がONになるとコレクタ電圧が下がり、結合コンデンサを通じて相手のベースへ変化が伝わる。
相手がOFFとなった後、抵抗を通じてそのベース側の電圧が徐々に回復する。
切替え条件に達すると役割が逆転する。コンデンサの電圧は瞬間的には変わりにくいので、コレクタの急な変化とベースの指数変化を対にして観察する。
\subsection{離散単安定回路}
無入力時には決まった安定状態にいる。トリガで一時的に逆の状態になり、RCの変化がしきい値に達すると戻る。
微分用・スピードアップ用・時間を決める容量は、同じ「コンデンサ」でも役割が異なる。
資料図で各容量の接続先を追い、どの波形のどの部分に効くか対応させる。
\clearpage
\subsection{555非安定・単安定}
非安定では出力HIGH中に$R_A+R_B$を通って充電する。2/3電源に達すると出力LOWとなり、放電端子から$R_B$を通って放電する。1/3電源で再び充電へ移る。
単安定ではトリガにより充電を始め、容量が2/3電源に達したところで出力が元へ戻る。非安定の1/3からの充電と、単安定の0からの充電は時間式が異なる。

\section{RC充放電と状態}\label{sec:auto3}
時定数$\tau=RC$に対して、コンデンサ電圧は
\[
 V_C(t)=V_\infty+(V_0-V_\infty)e^{-t/(RC)}
\]
と変化する。しきい値$V_{\rm th}$までの時間は
$t=-RC\ln[(V_{\rm th}-V_\infty)/(V_0-V_\infty)]$。
非安定には安定状態がなく、単安定には一つ、双安定には二つある。
\section{555の時間計算}\label{sec:auto4}
555はコンデンサ電圧が$2V_{CC}/3$に達すると放電し、$V_{CC}/3$で再び充電を始める。
充放電で電圧差が半分になる時間は$RC\ln2$である。
\[
 t_H=0.693(R_A+R_B)C,\quad t_L=0.693R_BC,\quad
 D=\frac{R_A+R_B}{R_A+2R_B}.
\]
\textbf{例題：}$R_A=5\,\mathrm{k}\Omega,R_B=3\,\mathrm{k}\Omega,C=0.1\,\mu\mathrm{F}$なら、
$t_H=0.554\,\mathrm{ms},t_L=0.208\,\mathrm{ms}$、
$f\simeq1.31\,\mathrm{kHz},D\simeq72.7\,\%$。
単安定は0から$2V_{CC}/3$まで充電するため$t=\ln3R_AC\simeq1.099R_AC$。
$5.6\,\mathrm{k}\Omega,0.1\,\mu\mathrm{F}$なら幅は約$0.616\,\mathrm{ms}$。
\section{離散回路を読む}\label{sec:auto5}
コレクタ電圧が低い側のトランジスタは通常ON、高い側はOFFとなる。ベース電圧、結合コンデンサ、微分入力を追って次の状態に移る理由を確認する。
\writing{24mm}
\clearpage
\section{測定で区別する量}\label{sec:auto6}
周期$T$は同じ向きの連続するエッジ間隔、パルス幅は一つの状態の持続時間である。
デューティ比は高レベル時間$/T$であり、入力と出力で同じとは限らない。
単安定では入力のパルス幅より出力が長くなる場合もあるが、再トリガや過長トリガの扱いは回路条件に依存する。
\note{負電源を含む離散回路と5 Vの555回路を混同しない。配線変更は電源OFFで行い、電解コンデンサの極性を確認する。}
\section{測定からレポートへ：時間と状態を対応させる}\label{sec:auto7}
\subsection{波形を同じ時間軸で読む}
出力だけでなく、ベースまたはコンデンサ電圧を同時に描く。
出力が切り替わる縦線を引き、その時に内部電圧がどのしきい値に達したか説明する。
測定を一度に行えない場合は、共通トリガと同じ時間/divを使って対応を保つ。
\subsection{周期・幅を混ぜない計算例}
仮の555波形でHIGHが0.56 ms、LOWが0.21 msなら、
\[
 T=0.77\,\mathrm{ms},\qquad f=1/T\simeq1.30\,\mathrm{kHz},\qquad D=0.56/0.77\simeq72.7\,\%.
\]
理論値は実際の抵抗と容量で計算してから比べる。周期と幅の誤差は別々に求める。
単安定には自動繰返しの周期式を使わず、出力パルスの幅を比較する。
\subsection{指数式を解くと対数が現れる}
555非安定の充電は、最終値$V_{CC}$との差が$2V_{CC}/3$から$V_{CC}/3$へ半分になるので$t=RC\ln2$。
単安定の0から2/3電源への充電では、最終値との差が1/3になるので$t=RC\ln3$。
係数0.69や1.1を暗記するだけでなく、開始値・最終値・しきい値を先に書く。

\clearpage
\section{確認問題}\label{sec:auto8}
\begin{enumerate}
\item $R_B=75\,\mathrm{k}\Omega$で離散回路の$T_1=1\,\mathrm{ms},T_2=0.5\,\mathrm{ms}$となる$C_1,C_2$を求める（係数0.69）。\answerline
\item 555の$5\,\mathrm{V}$電源における上下しきい値を求める。\answerline
\item 単安定のRだけを2倍にしたら出力幅はどうなるか。\answerline
\item なぜ充放電時間の式には$\ln2$や$\ln3$が現れるか。\answerline
\end{enumerate}
\section*{解答の目安}\label{sec:auto9}
(1) 約$19.3\,\mathrm{nF},9.66\,\mathrm{nF}$。
(2) 約$3.33\,\mathrm{V},1.67\,\mathrm{V}$。
(3) 理想近似では2倍。
(4) 指数関数の充放電を、所定のしきい値に達する時間について解くため。
\section*{補足資料}\label{sec:auto10}
Texas Instruments, \textit{xx555 Precision Timers}：
\url{https://www.ti.com/lit/ds/symlink/ne555.pdf}（参照日：2026-10-09）。
\prelabclosing
\end{document}
